\begin{statement}{E.29}{Proved}{All retained fibres survive the complete $A_5$ Ferrari test}
Let $D_{23}$ be the actual joint three-strand source quotient of
E.24, and let $T_{23}=T_{2,6}\times T_{3,6}$ be its actual
four-strand target. We retain all $432$ multiplier-one points
of E.24, including both binary corrections and every point
of each ternary plane.

There is an explicit finite Ferrari quotient
$\pi_A:B_4\twoheadrightarrow T_A$ whose free Ferrari kernel is
$A_5^{19}$. Its associated six-evaluation, six-coset induced
homomorphism is $\rho_A:B=C_2*C_3\twoheadrightarrow G_A$.
Every one of the $432$ points has a correction in $N'$ that
preserves its complete $D_{23}$ image and satisfies both
hexagons and the full pentagon in the common refinement with
$G_A$ and $T_A$. The resulting three-strand map and the map
onto the actual joint four-strand image are surjective.
In particular, all $24$ retained CM fibres, including all
$12$ extra fibres, survive this nonsolvable finite test.
\end{statement}
\begin{proof}
We give the target, the complete correction image, the
arithmetic comparison used to select a point in that image,
and the joint-generation argument separately.

\paragraph{The exact target and induced source.}
Enumerate the $2280$ ordered generating pairs of $A_5$.
Two such pairs give the same epimorphism kernel of the
free group $\langle\xi_1,\xi_2\rangle$ precisely when an
automorphism of $A_5$ relates them. Here
$\operatorname{Aut}(A_5)=S_5$, with its conjugation action.
For completeness, an automorphism permutes the five Sylow
$2$-subgroups, each isomorphic to $C_2^2$. Its action on this
set is faithful: an automorphism fixing the five subgroups
gives the same permutation of them for each $g$ and its
image, and the conjugation action of $A_5$ on these five
subgroups is faithful. Thus the automorphism group has
order at most $120$, while the $120$ conjugations by $S_5$
are distinct. This proves the equality used in the enumeration.

There are $19$ epimorphism kernels. Choose the lexicographic
representative $(g_j,h_j)$ of each orbit and retain all
$19$ evaluations. The Nielsen substitutions
\[
 a:(g,h)\longmapsto(g,hg),\qquad
 b:(g,h)\longmapsto(gh^{-1},h)
 \tag{E.29.1}
\]
permute these representatives, with a uniquely recorded
conjugating element of $S_5$ in each factor. The action on
the kernels has two orbits, of sizes $9$ and $10$.
The free evaluation image is $A_5^{19}$: a subdirect
product of simple groups can identify two factors only
through an isomorphism between the two corresponding
epimorphism kernels, which have already been separated.
The program also computes the order $60^{19}$ independently
in the resulting faithful permutation representation.

An element of the Ferrari semidirect product is recorded
as $(p,L,R)$, where $p$ permutes the $19$ factors, $R_j=c_j$
is the Nielsen conjugator, and $L_j=k_jc_j$ includes the
free-kernel coordinate. The multiplication used is
\[
 (p,L,R)(q,L',R')=
 \bigl(q\circ p,(L_jL'_{p(j)})_j,(R_jR'_{p(j)})_j\bigr).
 \tag{E.29.2}
\]
Two five-point blocks in each factor, acted on by $L_j$
and $R_j$, give a faithful permutation representation of
degree $190$. The literal images of the three braid
generators are $(1,a),(1,b),(\xi_1,a)$ in the Ferrari
description. The program checks every braid relation,
the literal words for $\xi_1,\xi_2$, and all four
Nielsen conjugation identities. The resulting orders are
\[
 \begin{split}
 |T_A|&=2^{60}3^{35}5^{31}7^2,\\
 |\langle a,b\rangle|&=2^{22}3^{16}5^{12}7^2.
 \end{split}
 \tag{E.29.3}
\]
All six evaluation pairs from E.17 are noncommuting in
this target.

For each of those six actual maps $\theta_i:F\to T_A$,
induce through the six cosets of $F\lhd B$. Explicitly,
for a section $q\mapsto b_q$ of $B/F$, the component at
$(q,i)$ of $\rho_A(w)$ is
\[
 \theta_i(b_qwb_{qw}^{-1}),
 \tag{E.29.4}
\]
together with the actual permutation $q\mapsto qw$.
This defines a homomorphism of discrete groups, extending
continuously to their profinite completions. Its kernel
is normal in $B$ and lies in $F$. On $F$ the coordinate
at the identity coset is exactly $\theta_i$; thus every
comparison with the six substitutions is a specified
coordinate homomorphism from $\rho_A(F)$.

The faithful induced permutation action has degree $6840$.
The executed compression to degree $1980$ deletes only
duplicate transitive $B$-sets: every deletion has a saved
bijective intertwiner for both $s$ and $r$. Both actions
give the identical image order
\[
 |G_A|=2^{38}3^{26}5^{21}7^2.
 \tag{E.29.5}
\]

\paragraph{The whole correction image.}
Write
\[
 P=\rho_A(N),\quad C=\rho_A(N'),\quad
 K=\rho_A(K_6),\quad E=\rho_A(F'),
 \qquad K_6=N'N^6.
\]
The exact permutation-group computations give
\[
 \begin{array}{c|c}
 \text{group}&\text{order or quotient}\\\hline
 P&2^{35}3^{24}5^{21}7^2\\
 C&2^{35}3^{23}5^{19}7^2\\
 K&2^{35}3^{23}5^{21}7^2\\
 P/C&75\\
 K/C&25\\
 G_A/P&72\\
 G_A/K&216\\
 E/C&4
 \end{array}
 \tag{E.29.6}
\]
Moreover,
\[
 C'=C,\qquad \rho_A(K_6\cap F')=C.
 \tag{E.29.7}
\]
These are complete subgroup calculations. The $97$
literal Schreier generators of $N$ give $P$; six of them
generate the computed image. The derived subgroup of
that group gives $C$. To compute $K_6\cap F'$ we retain
all $95$ products $\prod_i n_i^{6v_i}$, with $v$ running
through the primitive integral kernel of
$N^{\mathrm{ab}}\to F^{\mathrm{ab}}$ certified in E.19.
Together with $N'$ these normally generate
$K_6\cap F'$. All $95$ image products lie in $C$.
The full $97$ sixth powers instead enlarge $C$ by a
factor $25$, giving $K$. Stability under the actual
images of both $s$ and $r$ is checked. A separate
derived-subgroup computation on a reduced generating
set proves $C'=C$.

The modular comparison is also explicit. Map
\[
 s\longmapsto\begin{bmatrix}0&1\\-1&0\end{bmatrix},
 \qquad
 r\longmapsto\begin{bmatrix}0&-1\\1&-1\end{bmatrix}
 \quad\text{in }\operatorname{PSL}_2(\mathbb Z/6).
 \tag{E.29.8}
\]
Enumerate its $72$ elements and all $144$ Schreier
kernel generators for these two generators. Every
such word has its $\rho_A$ image in $P$. The two
preimages have index $72$, so they coincide. This
proves, with the displayed maps,
\[
 G_A/P\simeq\operatorname{PSL}_2(\mathbb Z/6)
       \simeq S_3\times A_4.
 \tag{E.29.9}
\]
Reducing (E.29.8) modulo $3$ defines $q_3:G_A\to A_4$.
For $u=r^{-1}s$, $v=sr^{-1}$ and $x=u^2,y=v^2$,
the two matrices are exactly
\[
 x\longmapsto\begin{bmatrix}1&2\\0&1\end{bmatrix},
 \qquad y\longmapsto\begin{bmatrix}1&0\\-2&1\end{bmatrix}.
 \tag{E.29.10}
\]
Their generated image is $A_4$, so $q_3(E)=A_4'=V_4$.
Since $C\subset P\subset\ker q_3$ and $|E/C|=4$,
\[
 \ker(q_3|_E)=C.
 \tag{E.29.11}
\]

\paragraph{Actual arithmetic points in the three required cosets.}
Nakamura--Schneps, Theorem~2.5
\cite{BCNakamuraSchneps2000}, gives for an actual
$\sigma\in G_{\mathbb Q}$ with $\chi(\sigma)=1$
\[
 q_3(f_\sigma)=M_{\kappa},\qquad
 M_\kappa=
 \begin{bmatrix}
 1&-8\kappa\\-8\kappa&1+64\kappa^2
 \end{bmatrix}\pmod{3},
 \quad \kappa=\rho_2(\sigma)\pmod3.
 \tag{E.29.12}
\]
The theorem's input matrices are exactly (E.29.10),
and its convention is
$x\mapsto x^\chi$, $y\mapsto f_\sigma^{-1}y^\chi f_\sigma$.
We use this theorem only for actual Galois elements.

All three values $\kappa=0,1,2$ occur with $\chi=1$
as an equality in $\widehat{\mathbb Z}^{\times}$.
Indeed, the splitting field
$L=\mathbb Q(\zeta_3,\sqrt[3]{2})$ has Galois group
$S_3$ and maximal abelian subextension $\mathbb Q(\zeta_3)$.
Since every finite subextension of
$\mathbb Q(\mu_\infty)/\mathbb Q$ is abelian,
\[
 L\cap\mathbb Q(\mu_\infty)=\mathbb Q(\zeta_3).
\]
Consequently $G_{\mathbb Q(\mu_\infty)}$ maps onto
$\operatorname{Gal}(L/\mathbb Q(\zeta_3))=C_3$,
and this restriction is precisely the Kummer coordinate
$\rho_2\bmod3$ used in (E.29.12).

Direct substitution of the $24$ retained E.17 words into
(E.29.10) gives the following complete index table:
\[
 \begin{array}{c|l}
 \kappa&\text{indices}\\\hline
 0&0,1,2,3,16,17,18,19\\
 1&4,5,6,7,20,21,22,23\\
 2&8,9,10,11,12,13,14,15
 \end{array}
 \tag{E.29.13}
\]
Every subsequent E.24 correction lies in $K_6$ and
therefore has trivial $q_3$ image. Thus the table applies
to all $432$ multiplier-one points, without discarding
any of their parameters.

Fix one of those literal words $f_0\in F'$ and choose
an actual $\sigma$ with the matching value in
(E.29.13). By (E.29.11),
\[
 c=\rho_A(f_0)^{-1}\rho_A(f_\sigma)\in C.
 \tag{E.29.14}
\]
Let $S$ be the image of $N'$ in $D_{23}$. It is
solvable, because $D_{23}$ is an extension of the
solvable group $Q$ by the explicitly constructed
nilpotent $2$- and $3$-groups. The joint image of $N'$
in $C\times S$ is the full product: Goursat's lemma
would otherwise give a nontrivial common quotient,
which would be both perfect and solvable. Hence
there is a discrete $k\in N'$ with images $(c,1)$.
For $f=f_0k$, we now have the exact equalities
\[
 f\bmod D_{23}=f_0\bmod D_{23},\qquad
 \rho_A(f)=\rho_A(f_\sigma),\qquad f\in F'.
 \tag{E.29.15}
\]
The same argument applies to the already constructed
solvable boundary refinements: at every finite level
the joint image is a product, and compactness supplies
one compatible correction in $\overline{N'}$ for their
inverse limit. This preserves their component maps and
equations. The simultaneous source-generation assertion
here is (E.29.16), for $D_{23}$; the corresponding
four-strand assertion is (E.29.18). We make no further
joint-generation assertion for a boundary refinement
whose order admits an additional common quotient of
$G_A/C$. This preservation argument also makes no
nonemptiness assertion over arbitrary profinite quotients.

The equality in $G_A$ matches all $36$ induced coordinates,
so it matches the six actual $T_A$ values of $f$ with
those of $f_\sigma$. In particular, the full pentagon,
the braid relations, and the full-twist identities in
$T_A$ follow from the actual Galois automorphism.
The source hexagons follow by the same equality in
$G_A$. All these equations remain true in $D_{23}$ and
$T_{23}$ by the first equality in (E.29.15).

\paragraph{Surjectivity in the common refinements.}
Put $H=G_A/K$. The actual $B$ image in
$D_{23}\times G_A$ is
\[
 D_{23}\mathbin{\times_H}G_A.
 \tag{E.29.16}
\]
Here $H$ is a quotient of $G_6$, with its specified
map from $B$, so both factors map onto it. To prove
the equality, observe that every solvable quotient of
$G_A$ kills the perfect group $C$. Every quotient of
order divisible only by $2$ and $3$ also kills the
normal subgroup $K/C$ of order $25$. Thus any common
quotient of $D_{23}$ and $G_A$ factors through $H$.
Their actual joint image already has the common
$H$ coordinate, so Goursat's lemma proves (E.29.16).
The corrected shadow surjects onto $D_{23}$ by E.24,
and onto $G_A$ because it has the same generator
images as $\rho_A$ composed with an actual Galois
automorphism of $\widehat B$. Its joint image therefore
equals (E.29.16), by the same common-quotient argument.

There is an equally complete four-strand comparison.
A separate computation gives
\[
 [T_A:T_A']=2,\qquad T_A''=T_A'.
 \tag{E.29.17}
\]
All three braid generators have the same nontrivial
image in $T_A/T_A'=C_2$. The exact derived series is
computed in the faithful $190$-point action; the
subgroup orders and the generator-ratio memberships
are retained in the certificate. Thus every solvable
quotient of $T_A$ factors through this parity map.
The group $T_{23}$ is solvable by its explicit
$2$- and $3$-group extensions and has the same parity
map through its binary $S_4$ quotient. Therefore the
actual joint target is
\[
 T_{23}\mathbin{\times_{C_2}}T_A.
 \tag{E.29.18}
\]
The corrected three braid generators surject onto
$T_{23}$ by E.24 and onto $T_A$ because they are the
image of the Galois braid automorphism in that factor.
They preserve the common parity coordinate. Goursat's
lemma and (E.29.17) then prove surjectivity onto
(E.29.18). Its kernel in $B_4$ is normal and lies in
$PB_4$, since the standard $S_4$ map already factors
through $T_{23}$. The coordinate maps (E.29.4) give
the required comparisons for every pentagon
substitution.

For the braid-action convention, we use
Nakamura--Schneps (1.1), which gives
$\sigma_i\mapsto
 f(y_i,\sigma_i^2)^{-1}\sigma_i^\lambda
 f(y_i,\sigma_i^2)$.
For $i=3$, our sixth pair uses
$\Omega_3=(ab)^3=a^2(ba^2b)$ instead of $y_3=ba^2b$.
The element $a^2$ commutes with both $y_3$ and $c^2$;
as $f\in F'$, inserting that commuting factor does
not change the evaluation. Hence the exact six-pair
action used here is the stated Galois action.

The executed certificates are
\texttt{de\_ferrari\_a5\_model.py},
\texttt{de\_ferrari\_a5\_components.py},
\texttt{de\_ferrari\_a5\_source.py},
\texttt{de\_ferrari\_a5\_source\_compress.py},
\texttt{de\_ferrari\_a5\_corrections.py},
\texttt{de\_ferrari\_a5\_compare.py}, and
\texttt{de\_ferrari\_a5\_solvable.py}, with their
outputs and exact permutation arrays. The first new
correction is supplied by (E.29.14)--(E.29.15); no
uncomputed literal Galois word is claimed as a saved
certificate. The finite group and fibre identifications,
and every numerical assertion used in that existence
proof, were computed and retained.
\end{proof}
